% ЗГЕНЕРОВАНО scripts/build_format_v2_chapter.py — не редагувати вручну.
% Джерело: editorial/cards/6.*.json, format-v2/rozdil-6.json, paraleli.json.

\naukchapter{6}{4}{Глава 6. Горщик і простір}
\naukintro{Шоста глава тримається на рефрені, довшому за рефрен пʼятої: чотири вірші закінчуються тими самими словами про відкидання, взяття й розчинення. Кожна строфа дає новий образ — простір і горщик, океан і хвиля, перламутр і срібло, — і після кожного повторює, що з цим знанням нічого не треба робити. Це глава про те, що знання не є дією.}
\vsec{6.1}
\vdeva{आकाशवद् अनन्तो ऽहं घटवत् प्रकृतं जगत् । इति ज्ञानं तथैतस्य न त्यागो न ग्रहो लयः ॥}
\vtranslit{акашавад ананто ’хам ґгатават пракрітам джаґат | іті джнянам татхайтася на тяґо на ґрахо лаях ||}
\vperek{Я безмежний, як простір, а світ — як горщик. Таке знання — і в нього немає ні відкидання, ні взяття, ні розчинення.}
\vcom{Построкове зіставлення: редакторська чернетка, не рецензовано.}
\vcom{Рядок а}
\vdeva{आकाशवद् अनन्तो ऽहं घटवत् प्रकृतं जगत्}
\vtranslit{акашавад ананто ’хам ґгатават пракрітам джаґат}
\vperek{Я безмежний, як простір, а світ — як горщик.}
\vcom{ākāśavad і ghaṭavat — порівняння «як простір» та «як горщик». ananto 'ham стосується безмежного «я», тоді як prakṛtaṃ jagat — створеного світу. Це метафоричне зіставлення, а не буквальний опис простору.}
\vcom{Рядок б}
\vdeva{इति ज्ञानं तथैतस्य न त्यागो न ग्रहो लयः}
\vtranslit{іті джнянам татхайтася на тяґо на ґрахо лаях}
\vperek{Таке знання — і в нього немає ні відкидання, ні взяття, ні розчинення.}
\vcom{iti jñānaṃ — «таке знання»; na tyāgo na graho layaḥ заперечує відкидання, взяття і розчинення. Сенс рефрену — відсутність потреби щось робити з таким знанням, а не вимога знищити світ.}
\vcom{Образ простору й горщика працює точніше, ніж здається: простір не вміщується в горщик і не обмежується його стінками — горщик лише окреслює частину простору, не додаючи й не забираючи нічого. Тому рефрен і каже: немає ні tyāga (відкидання), ні graha (взяття), ні навіть laya (розчинення) — усі три дії припускають, що щось справді додалося або зникло. Слово prakṛtam («зроблене») підкреслює: горщик — виріб, а простір — ні. РЕФРЕН ГЛАВИ. Усі чотири вірші глави 6 закінчуються однаково: iti jñānaṃ tathaitasya na tyāgo na graho layaḥ («таке знання — і в нього немає ні відкидання, ні взяття, ні розчинення»). Зверни увагу: рефрен глави 5 вимагав laya («йди до розчинення»), а рефрен глави 6 заперечує навіть laya. Це не суперечність, а наступний крок: спершу розчинення як те, до чого йдуть, потім — знання, у якому немає кому йти й нічого розчиняти. Главу слід читати як одне ціле.}
\vpar{Українська паралель (ПОТРЕБУЄ ПЕРЕВІРКИ)}{українська філософська традиція, не визначено.}{ПОТРЕБУЄ ПЕРЕВІРКИ: образ простору, що не ділиться стінками посуду, найчастіше приходить до українського читача через Сковороду; саме місце, де він говорить про посуд і простір, ми за джерелом не звірили й тому не називаємо.}
\vfig{ilustratsii/6-1-space.pdf}{Простір і горщик: стінки ділять посуд, а не простір}{Схема наша, не історична ілюстрація.}

\vsec{6.2}
\vdeva{महोदधिर् इवाहं स प्रपञ्चो वीचिसंनिभः । इति ज्ञानं तथैतस्य न त्यागो न ग्रहो लयः ॥}
\vtranslit{маходадгір івахам са прапаньчо вічісамнібгах | іті джнянам татхайтася на тяґо на ґрахо лаях ||}
\vperek{Я — як великий океан, а цей плин світу — як хвилі. Таке знання — і в нього немає ні відкидання, ні взяття, ні розчинення.}
\vcom{Построкове зіставлення: редакторська чернетка, не рецензовано.}
\vcom{Рядок а}
\vdeva{महोदधिर् इवाहं स प्रपञ्चो वीचिसंनिभः}
\vtranslit{маходадгір івахам са прапаньчо вічісамнібгах}
\vperek{Я — як великий океан, а цей плин світу — як хвилі.}
\vcom{mahodadhir ivāhaṃ — «я, мов великий океан»; prapañcaḥ — розгортання світу; vīcisaṃnibhaḥ — «подібний до хвиль». Хвиля тут є образом явищ, а не незалежною від океану сутністю.}
\vcom{Рядок б}
\vdeva{इति ज्ञानं तथैतस्य न त्यागो न ग्रहो लयः}
\vtranslit{іті джнянам татхайтася на тяґо на ґрахо лаях}
\vperek{Таке знання — і в нього немає ні відкидання, ні взяття, ні розчинення.}
\vcom{Спільний рефрен iti jñānaṃ ... na tyāgo na graho layaḥ заперечує потребу брати, відкидати або розчиняти явища. Порівняння океану з хвилями не є твердженням про фізичну будову світу.}
\vcom{prapañca — важке слово: буквально «розгортання, розпросторення», а в ужитку «різноманіття світу» з відтінком заплутаності. Тут воно вжите без засудження: хвилі не є вадою океану. Порівняй 2.23–2.25, де той самий образ розгорнутий у три кроки; тут він стиснутий до одного порівняння, бо мета інша — не показати, як хвилі виникають, а показати, що з ними не треба нічого робити. РЕФРЕН ГЛАВИ. Усі чотири вірші глави 6 закінчуються однаково: iti jñānaṃ tathaitasya na tyāgo na graho layaḥ («таке знання — і в нього немає ні відкидання, ні взяття, ні розчинення»). Зверни увагу: рефрен глави 5 вимагав laya («йди до розчинення»), а рефрен глави 6 заперечує навіть laya. Це не суперечність, а наступний крок: спершу розчинення як те, до чого йдуть, потім — знання, у якому немає кому йти й нічого розчиняти. Главу слід читати як одне ціле.}
\vpar{Українська паралель (наша інтерпретація)}{Григорій Сковорода, «Вдячний Еродій» (трактат).}{Наша паралель: у Сковороді море як образ цілого, у якому частка не має власного окремого життя, близьке до цього вірша. Наше порівняння — типологічне, не твердження про джерело.}
\vfig{ilustratsii/6-2-motion.pdf}{Хвилі й океан: у знанні немає ні взяття, ні відкидання}{Схема наша, не історична ілюстрація.}

\vsec{6.3}
\vdeva{अहं स शुक्तिसङ्काशो रूप्यवद् विश्वकल्पना । इति ज्ञानं तथैतस्य न त्यागो न ग्रहो लयः ॥}
\vtranslit{ахам са шуктісанкашо руп'явад вішвакалпана | іті джнянам татхайтася на тяґо на ґрахо лаях ||}
\vperek{Я — як перламутр, а уявлення світу — як срібло в ньому. Таке знання — і в нього немає ні відкидання, ні взяття, ні розчинення.}
\vcom{Построкове зіставлення: редакторська чернетка, не рецензовано.}
\vcom{Рядок а}
\vdeva{अहं स शुक्तिसङ्काशो रूप्यवद् विश्वकल्पना}
\vtranslit{ахам са шуктісанкашо руп'явад вішвакалпана}
\vperek{Я — як перламутр, а уявлення світу — як срібло в ньому.}
\vcom{śuktisaṅkāśaḥ — «подібний до перламутру»; rūpyavat — «як срібло»; viśvakalpanā — уявлення світу. Порівняння показує різницю між тим, що бачать, і тим, що помилково приписують побаченому.}
\vcom{Рядок б}
\vdeva{इति ज्ञानं तथैतस्य न त्यागो न ग्रहो लयः}
\vtranslit{іті джнянам татхайтася на тяґо на ґрахо лаях}
\vperek{Таке знання — і в нього немає ні відкидання, ні взяття, ні розчинення.}
\vcom{Рефрен na tyāgo na graho layaḥ заперечує відкидання, взяття й розчинення: у метафорі не потрібно знищувати неіснуюче «срібло». Це пояснювальна аналогія хибного сприйняття, не рукописний доказ.}
\vcom{Третє й останнє з трьох порівнянь глави (простір/горщик, океан/хвилі, перламутр/срібло), і воно єдине, де друга річ не існує зовсім: горщик і хвиля є, срібла немає. Тому тут і стоїть kalpanā («уявлення, конструкція») — не «світ», а саме уявлення світу. Зверни увагу: перламутр не заперечується і не змінюється — зникає тільки срібло, а разом із ним і питання, що з ним робити. Це і є сенс рефрену. РЕФРЕН ГЛАВИ. Усі чотири вірші глави 6 закінчуються однаково: iti jñānaṃ tathaitasya na tyāgo na graho layaḥ («таке знання — і в нього немає ні відкидання, ні взяття, ні розчинення»). Зверни увагу: рефрен глави 5 вимагав laya («йди до розчинення»), а рефрен глави 6 заперечує навіть laya. Це не суперечність, а наступний крок: спершу розчинення як те, до чого йдуть, потім — знання, у якому немає кому йти й нічого розчиняти. Главу слід читати як одне ціле.}
\vpar{Українська паралель (ПОТРЕБУЄ ПЕРЕВІРКИ)}{українська барокова емблематика, не визначено.}{ПОТРЕБУЄ ПЕРЕВІРКИ: пара перламутр / срібло як образ оманливого уявлення знайома бароковій емблематиці, але конкретного твору ми не перевірили.}
\vfig{ilustratsii/6-3-illusion.pdf}{Перламутр і срібло: уявлення світу}{Схема наша, не історична ілюстрація.}

\vsec{6.4}
\vdeva{अहं वा सर्वभूतेषु सर्वभूतान्य् अथो मयि । इति ज्ञानं तथैतस्य न त्यागो न ग्रहो लयः ॥}
\vtranslit{ахам ва сарвабгутешу сарвабгутанй атхо маї | іті джнянам татхайтася на тяґо на ґрахо лаях ||}
\vperek{Я — у всіх істотах, а всі істоти — у мені. Таке знання — і в нього немає ні відкидання, ні взяття, ні розчинення.}
\vcom{Построкове зіставлення: редакторська чернетка, не рецензовано.}
\vcom{Рядок а}
\vdeva{अहं वा सर्वभूतेषु सर्वभूतान्य् अथो मयि}
\vtranslit{ахам ва сарвабгутешу сарвабгутанй атхо маї}
\vperek{Я — у всіх істотах, а всі істоти — у мені.}
\vcom{sarvabhūteṣu — «у всіх істотах» (місцевий множини), sarvabhūtāni — «усі істоти», mayi — «у мені». atho поєднує два напрями взаємного перебування.}
\vcom{Рядок б}
\vdeva{इति ज्ञानं तथैतस्य न त्यागो न ग्रहो लयः}
\vtranslit{іті джнянам татхайтася на тяґо на ґрахо лаях}
\vperek{Таке знання — і в нього немає ні відкидання, ні взяття, ні розчинення.}
\vcom{Підсумкове iti jñānaṃ ... na tyāgo na graho layaḥ повторює заперечення відкидання, взяття й розчинення. Цей рефрен завершує главу після трьох образних порівнянь.}
\vcom{Двозмінна формула та сама, що в 3.5, і тут вона завершує главу. Після трьох порівнянь (простір, океан, перламутр) висновок поданий без образу взагалі — лише відношення. atho («а також») з'єднує два напрямки так, що вони не зливаються в один: я в істотах і істоти в мені — не те саме твердження, прочитане двічі, а два боки одного відношення. Саме тому рефрен знову заперечує tyāga і graha: у такому відношенні немає нічого, що можна було б узяти або відкинути. РЕФРЕН ГЛАВИ. Усі чотири вірші глави 6 закінчуються однаково: iti jñānaṃ tathaitasya na tyāgo na graho layaḥ («таке знання — і в нього немає ні відкидання, ні взяття, ні розчинення»). Зверни увагу: рефрен глави 5 вимагав laya («йди до розчинення»), а рефрен глави 6 заперечує навіть laya. Це не суперечність, а наступний крок: спершу розчинення як те, до чого йдуть, потім — знання, у якому немає кому йти й нічого розчиняти. Главу слід читати як одне ціле.}
\vpar{Українська паралель (наша інтерпретація)}{Павло Тичина, «Плуг» (1920).}{Наша паралель: у Тичини взаємне перебування всього в усьому теж стає способом бачити світ, а не твердженням про будову речей. Наше порівняння — за образом.}
\vfig{ilustratsii/6-4-witness.pdf}{Взаємне перебування: я в усіх істотах, усі істоти в мені}{Схема наша, не історична ілюстрація.}

